\documentclass{exam-zh}
\usepackage{edu-practice}

\examsetup{
  question/show-answer = true,
  fillin/show-answer = true,
  solution/show-solution = show-stay,
}

% ---------- 教师版解答样式 ----------
\newtcolorbox{solutionblock}{
  enhanced, breakable,
  colback=edu-blue-bg,
  colframe=edu-blue!25,
  boxrule=0pt,
  borderline west={2pt}{0pt}{edu-blue!50},
  arc=0pt,
  left=3mm, right=3mm, top=2mm, bottom=2mm,
  before skip=2mm,
  after skip=3mm,
  fontupper=\small,
}

% 解答步骤编号
\newcounter{solstep}
\newcommand{\solstep}[2]{%
  \stepcounter{solstep}%
  \par\medskip
  \noindent\hangindent=6mm
  {\color{edu-blue}\bfseries\textcircled{\scriptsize\thesolstep}\enspace #1}\enspace #2\par
}

\begin{document}

\section*{线段与锐角三角比双向互求练习}

\section*{一、常见勾股数：先把方向做对}


\begin{question}[points=4]
  在 $\operatorname{Rt}\triangle ABC$ 中，$\angle B=90^\circ$，$AB=3$，$BC=4$，则 $\sin C=$（\quad）
  
  \begin{choices}
    \item $\dfrac34$
    \item $\dfrac35$
    \item $\dfrac45$
    \item $\dfrac43$
  \end{choices}
  \paren[B]
\end{question}



\begin{question}[points=4]
  在 $\operatorname{Rt}\triangle ABC$ 中，$\angle B=90^\circ$，$AB=5$，$BC=12$，则 $\tan C=$\fillin。
  
  \paren[$\dfrac5{12}$]
\end{question}


\needspace{8\baselineskip}

\begin{problem}[points=4]
    在 $\operatorname{Rt}\triangle ABC$ 中，$\angle B=90^\circ$，
$\sin C=\dfrac7{25}$，$BC=24$，求 $AC$。

    
  \answerarea[28mm]
  \setcounter{solstep}{0}
  \begin{solutionblock}
    \textbf{答案：}$AC=25$。\par\medskip
    \solstep{还原份数}{$\sin C=AB/AC=7/25$，故 $AB:BC:AC=7:24:25$。}
    \solstep{由已知边定倍数}{$BC=24$，缩放倍数为 $1$，所以 $AC=25$。}
    
  \end{solutionblock}
\end{problem}


\section*{二、常见勾股数：交换已知与所求}


\begin{question}[points=4]
  在 $\operatorname{Rt}\triangle ABC$ 中，$\angle B=90^\circ$，$\cos C=\dfrac{15}{17}$，$AC=17$，则 $AB=$\fillin。
  
  \paren[$8$]
\end{question}


\needspace{8\baselineskip}

\begin{problem}[points=4]
    在 $\operatorname{Rt}\triangle ABC$ 中，$\angle B=90^\circ$，
$\tan C=\dfrac34$，$AC=10$。求 $AB$ 和 $BC$。

    
  \answerarea[32mm]
  \setcounter{solstep}{0}
  \begin{solutionblock}
    \textbf{答案：}$AB=6$，$BC=8$。\par\medskip
    \solstep{还原三边份数}{$\tan C=AB/BC=3/4$，故 $AB:BC:AC=3:4:5$。}
    \solstep{同倍缩放}{$AC=10=5\times2$，所以 $AB=3\times2=6$，$BC=4\times2=8$。}
    
  \end{solutionblock}
\end{problem}



\begin{question}[points=4]
  在 $\operatorname{Rt}\triangle ABC$ 中，$\angle B=90^\circ$，$AB=\dfrac{12}{7}$，$BC=\dfrac{16}{7}$，则 $\sin C=$\fillin。
  
  \paren[$\dfrac35$]
\end{question}


\section*{三、分数倍勾股数：从边求比、从比求边}


\begin{question}[points=4]
  在 $\operatorname{Rt}\triangle ABC$ 中，$\angle B=90^\circ$，$AB=\dfrac52$，$BC=6$，则 $\cos C=$（\quad）
  
  \begin{choices}
    \item $\dfrac5{12}$
    \item $\dfrac5{13}$
    \item $\dfrac{12}{13}$
    \item $\dfrac{13}{12}$
  \end{choices}
  \paren[C]
\end{question}



\begin{question}[points=4]
  在 $\operatorname{Rt}\triangle ABC$ 中，$\angle B=90^\circ$，$\tan C=\dfrac34$，$AC=\dfrac52$，则 $AB=$\fillin。
  
  \paren[$\dfrac32$]
\end{question}


\needspace{8\baselineskip}

\begin{problem}[points=4]
    在 $\operatorname{Rt}\triangle ABC$ 中，$\angle B=90^\circ$，
$\cos C=\dfrac45$，$AC=4$，求 $BC$。

    
  \answerarea[28mm]
  \setcounter{solstep}{0}
  \begin{solutionblock}
    \textbf{答案：}$BC=\dfrac{16}{5}$。\par\medskip
    \solstep{读出边比}{$\cos C=BC/AC=4/5$。}
    \solstep{按斜边缩放}{$AC=4$，故 $BC=4\times\dfrac45=\dfrac{16}{5}$。}
    
  \end{solutionblock}
\end{problem}


\section*{四、分数倍勾股数：识别公因子}


\begin{question}[points=4]
  在 $\operatorname{Rt}\triangle ABC$ 中，$\angle B=90^\circ$，$AB=\dfrac{10}{3}$，$BC=8$，则 $\tan C=$\fillin。
  
  \paren[$\dfrac5{12}$]
\end{question}


\needspace{8\baselineskip}

\begin{problem}[points=4]
    在 $\operatorname{Rt}\triangle ABC$ 中，$\angle B=90^\circ$，
$\sin C=\dfrac35$，$BC=5$，求 $AC$。

    
  \answerarea[28mm]
  \setcounter{solstep}{0}
  \begin{solutionblock}
    \textbf{答案：}$AC=\dfrac{25}{4}$。\par\medskip
    \solstep{补全份数}{$\sin C=3/5$，故 $AB:BC:AC=3:4:5$。}
    \solstep{确定倍数}{$BC=5=4\times\dfrac54$，故 $AC=5\times\dfrac54=\dfrac{25}{4}$。}
    
  \end{solutionblock}
\end{problem}



\begin{question}[points=4]
  在 $\operatorname{Rt}\triangle ABC$ 中，$\angle B=90^\circ$，$AB=\dfrac{21}{5}$，$BC=\dfrac{28}{5}$，则 $\cot C=$（\quad）
  
  \begin{choices}
    \item $\dfrac34$
    \item $\dfrac43$
    \item $\dfrac35$
    \item $\dfrac45$
  \end{choices}
  \paren[B]
\end{question}


\section*{五、双向混合：不看题号，只看已知与所求}


\begin{question}[points=4]
  在 $\operatorname{Rt}\triangle ABC$ 中，$\angle B=90^\circ$，$\cos C=\dfrac45$，$AC=3$，则 $AB=$\fillin。
  
  \paren[$\dfrac95$]
\end{question}


\needspace{8\baselineskip}

\begin{problem}[points=4]
    在 $\operatorname{Rt}\triangle ABC$ 中，$\angle B=90^\circ$，
$\tan C=\dfrac34$，$AC=\dfrac{35}{4}$，求 $BC$。

    
  \answerarea[28mm]
  \setcounter{solstep}{0}
  \begin{solutionblock}
    \textbf{答案：}$BC=7$。\par\medskip
    \solstep{还原三边份数}{$\tan C=3/4$，故 $AB:BC:AC=3:4:5$。}
    \solstep{同倍缩放}{$AC=35/4=5\times7/4$，所以 $BC=4\times7/4=7$。}
    
  \end{solutionblock}
\end{problem}


\needspace{8\baselineskip}

\begin{problem}[points=4]
    在 $\operatorname{Rt}\triangle ABC$ 中，$\angle B=90^\circ$，
$AB=2$，$BC=\dfrac83$。求 $\sin C$ 和 $\cos C$。

    
  \answerarea[34mm]
  \setcounter{solstep}{0}
  \begin{solutionblock}
    \textbf{答案：}$\sin C=\dfrac35$，$\cos C=\dfrac45$。\par\medskip
    \solstep{化回熟悉份数}{$AB:BC=2:\dfrac83=3:4$，故 $AB:BC:AC=3:4:5$。}
    \solstep{分别列比}{$\sin C=AB/AC=3/5$，$\cos C=BC/AC=4/5$。}
    
  \end{solutionblock}
\end{problem}


\clearpage
\section*{答案}




\end{document}
